\section{Introduction}
\label{sec:introduction}

A change of coordinates can relate a prescribed sequence of matching
directions to a random one. For a permutation shuffle, we identify
the joint path laws under the change of coordinates. A conditional
calculation also needs to specify what
has been revealed when the next direction is sampled. We separate these
two questions. The resulting estimates describe the laws of large
ordered lists of cards, and several ways of retaining that information
lead to convergence of the whole permutation.

Our model is the Thorp shuffle on $n=2^d$ labeled cards. Split a deck
into two equal halves, pair corresponding positions, and interleave the
pairs, choosing their orders with independent fair coins. One such
operation is one \emph{physical shuffle}. Let $q_d$ be its law as a
permutation of positions, let $U_n$ be uniform on $S_n$, and let
$\mu*\nu$ denote the law of $XY$ for independent permutations with laws
$\mu,\nu$. We use
$\|\mu-\nu\|_{\TV}=\frac12\sum_g|\mu(g)-\nu(g)|$.
A card becomes uniform after $d$ physical shuffles. The cards share the
same switches, so the law of their joint positions requires a separate
argument.

Our first complete application uses the frame, marginal, degree, and
Fourier arguments proved in this paper.

\begin{theorem}\label{thm:main}
For $n=2^d$,
\[
 \|q_d^{*(32800d)}-U_n\|_{\TV}\longrightarrow0
 \qquad(d\longrightarrow\infty).
\]
The convergence is uniform over every deterministic initial ordering
of the deck.
\end{theorem}

The lower obstruction is elementary. After $t$ physical shuffles the
process has used $tn/2$ random bits, so its support has size at most
$2^{tn/2}$ and
\[
 \|q_d^{*t}-U_n\|_{\TV}\ge1-\frac{2^{tn/2}}{n!}.
\]
At distance $1/4$ this forces
$t\ge\lceil(2/n)\log_2(3n!/4)\rceil=2d-O(1)$.
Together with Theorem~\ref{thm:main}, this shows that the mixing time
has the optimal order $d=\log_2 n$ on dyadic decks.

The proof of Theorem~\ref{thm:main} is independent of the
deterministic-coordinate companion, which proves convergence after
$1600d$ physical shuffles as $d\to\infty$
\cite[Theorem~1.1]{optimal-thorp}, and of the abstract Fourier
transfers in \cite{partial-fourier}. Later we prove that, after $256d$
shuffles, the ordered image of every fixed list of at most $7n/8$ cards
has total-variation distance $o(1)$ from a uniform injection, uniformly
over the list as $d\to\infty$. The orbit transfer gives the $2048d$ application in
Theorem~\ref{shear:thm:2048}. Combining the same marginal estimate with
the Fourier companion's isotypic estimate and its inverse-degree sum gives
the $512d$ consequence in Corollary~\ref{shear:cor:512}. The orbit
argument bounds the operator norm, while the isotypic argument bounds
the Hilbert--Schmidt norm. Later sections also prove partial laws that
retain information about histories and changing input sets.

\subsection{Coordinate frames and the first proof}

Identify positions with $V=\mathbb F_2^d$. Undoing the deterministic
rotation in each shuffle leaves independent fair switches along
successive coordinate directions. A run through all $d$ directions is
a \emph{sweep}; it costs $d$ physical shuffles. We compare this process
with switches in directions $v_1,v_2,\ldots$ such that every $d$
consecutive vectors form a basis. For a fixed sequence whose first
basis is the same ordered coordinate basis, invertible changes of position frame
identify the joint path laws. At the terminal time, the endpoint laws
agree after a specified bijection of output positions. Hence every fixed input
list has the same distance from a uniform injection in the two
descriptions. A general initial frame also relabels the inputs; the
corresponding comparison preserves the average over uniformly chosen
input lists.

This path comparison is proved after fixing the complete direction
sequence. Its intermediate maps may use future directions. The
probability calculation therefore uses a separately generated process:
after the first basis, choose the next direction uniformly outside the
span of its $d-1$ predecessors, then choose fresh switch coins. The
observation field records the input list and any independent initial
sampling, the direction prefix, and the path prefixes
of the preceding cards in the list. In this field, revealing the next
direction does not change the conditional law of the hidden card's
current position.

That card has a conditional mass on the positions not occupied by its
predecessors. A switch averages two masses when both endpoints remain
available, and transports one mass when only one endpoint remains.
The squared deviation from uniform has an exact decrement. Averaging
a fresh direction connects every pair across the hyperplane spanned by
the preceding directions with the same probability. The expected decrement
is at least the current energy times the smaller available half-count,
divided by $n$.

For the first proof we sample an omitted block of $n/64$ labels
independently of the shuffle. Its
image is contained in every available set for the complementary list.
Once the shuffle history is fixed, that image is a uniform subset, so
both sides of a specified hyperplane contain many of its positions
except with exponentially small probability. This is an unconditional
imbalance bound after the path conditioning has been removed. Since
the energy is at most one, it can be combined with the decrement even
when the energy and the imbalance event are correlated. The terminal
conditional comparison keeps the complete direction data inside total
variation until the frame identity is applied. It yields one partition
into 64 blocks for which the sum of the complementary marginal errors
after $1025d$ shuffles tends to zero.

Each complementary image list identifies a left coset of the subgroup
that permutes the omitted block. A common trimming gives one nearby law
whose relevant coset masses are capped by a fixed multiple of their
uniform values. On restricting an
irreducible representation of dimension $D$ to the product of the 64
block groups, each occurring tensor type has a factor of degree at most
$D^{1/64}$. The orbit of one vector under
that factor has controlled dimension even when the factor occurs with
large multiplicity. The transfer proof first bounds this dimension
by a finite sum of squared subgroup degrees. A local estimate for
the inverse sixteenth powers then bounds that sum, and Schur averaging
turns the coset caps into Fourier contraction. Thirty-two independent
time blocks complete the proof, with the sign representation treated
separately.

Sections~\ref{sec:prelim}--\ref{sec:frames} establish the model and
the fixed-schedule path identity. Section~\ref{sec:marginals} proves
the marginal estimate. Sections~\ref{sec:lift}, \ref{sec:degrees},
and~\ref{sec:completion} complete Theorem~\ref{thm:main} using local
representation and Fourier arguments. A reader may stop there with a
complete proof.

\subsection{Further partial laws and their uses}

Section~\ref{shear:sec:directions} reuses the conditional mass flow with
different balance and disclosure statements. A preceding transverse
switch gives an imbalance bound for every fixed input list. Independent
shears provide an explicit realization of the fresh directions, and
the deferred-column calculation identifies the coefficients that have
not yet been revealed. Sampling labels offers other balance estimates;
an arbitrary initial frame also introduces an input relabeling that
disappears after averaging over lists. Thus an average over input lists
can enter through the frame comparison or through the balance estimate.
Lemma~\ref{shear:lem:delayed-symmetry} gives a separate contraction under an earlier
conditioning time for the deterministic coordinate schedule.

Section~\ref{frames:chapter} controls exceptional trajectories in three different ways.
For a chosen partition and chosen orders of its complements, descending
removal constructs one retained measure whose complementary marginal
probabilities are capped pointwise. A stopped argument uses nested balance
events so that the tuple error includes the imbalance probability only
once. A different event, defined by the
trajectory matchings, supports an entropy inequality for every
probability law concentrated on that event. This last quantifier allows
the Fourier transfer to condition on a rare event selected from a
representation coefficient.

Section~\ref{lifts:random-marginals} uses averaged input laws in three ways: selecting one partition
for an isotypic estimate, trimming equivariant kernels while controlling
their full sign multiplicity spaces, and averaging over overlapping
omitted sets. Sections~\ref{frames:char:section} and~\ref{sec:character}
use half of the deck. When
an independent pass follows a long mixing segment, separate caps on
preimage lists control positive trace tests. When the pass comes first,
the argument repairs the two image marginals and averages one uniform
subgroup coordinate at a time. The repaired coordinates may remain
dependent.

\subsection{History and related work}

The model arose in Thorp's study of nonrandom shuffling and the game
of Faro~\cite{thorp}. Aldous and Diaconis later described the paired
shuffle and reported an informal argument for an
$O((\log n)^2)$ threshold~\cite[Section~8b, Example~(12)]{aldous-diaconis1987}.
For dyadic decks, the description by independent switches along
successive hypercube coordinates is standard~\cite{morris44}. It makes
one-card uniformity transparent, while the dependence among card
trajectories remains the obstacle to convergence of the full law.

Morris obtained an $O(d^{44})$ bound for the full deck when $n=2^d$,
combining chameleon and evolving-set ideas~\cite{morris44}. Montenegro
and Tetali sharpened that analysis to $O(d^{29})$ using spectral-profile
and evolving-set estimates~\cite[Theorems~6.14--6.15]{montenegro-tetali2006}.
Morris subsequently used entropy estimates to prove
$O((\log n)^4)$ for every even deck size~\cite{morris4} and then
$O(d^3)$ for dyadic decks~\cite{morris3}. These times count individual
physical shuffles and concern the entire permutation.

Partial-card laws were studied both as shuffle marginals and as a tool
for enciphering small domains. Morris, Rogaway and Stegers gave explicit
bounds for ordered Thorp marginals~\cite[Theorem~1]{mrs}. For dyadic
decks, Czumaj and V\"ocking proved that any prescribed list of at most
$(1-\epsilon)n$ cards mixes in $O_\epsilon((\log n)^2)$ physical layers,
for each fixed $0<\epsilon<1$, using non-Markovian couplings in the
butterfly network~\cite{CzumajVocking2014}. Czumaj's later account
identifies this as $O(\log n)$ whole butterflies~\cite[Section~1.5]{Czumaj2015}.
That paper also obtains logarithmic-depth partial mixing for other
switching networks with suitable expansion~\cite[Theorems~3.4--3.6]{Czumaj2015}.

We adapt the conditional-energy approach in the swap-or-not analysis
of Hoang, Morris and Rogaway~\cite[Section~3, Theorem~3 and Lemma~4]{hmr2014}.
Their independent uniform group keys give an unconstrained sequence of
matchings; their proof contracts a quadratic discrepancy for the next card
and compares the joint law through expected conditional errors. Dai,
Hoang and Tessaro later sharpened the combination of those squared
discrepancies through a chi-squared inequality~\cite[Section~6]{dai-hoang-tessaro2017}.
Here the next direction is constrained by its recent predecessors.
The frame identity and the chronological conditional calculation
justify that rule and its hyperplane energy estimate.

Related work connects partial information to stronger permutation
guarantees in several ways. Ristenpart and Yilek's mix-and-cut construction,
and Morris and Rogaway's sometimes-recurse construction, use partial
or separator guarantees inside recursively modified shuffles
\cite{ristenpart-yilek2013,morris-rogaway2014}. Alon and Lovett relate
uniformity for every input-output pair of a group action to Fourier
matrices of the irreducible representations occurring in that action,
and repair sufficiently accurate approximate uniformity to exact
uniformity for the action~\cite[Theorem~1.4 and Proposition~3.2]{alon-lovett2013}.
Our passage to the full symmetric group controls its irreducible Fourier
matrices from complementary coset caps for selected blocks.
